\slajdid{07-s002}
\begin{frame}[label=07-s002]{Opis digitalnog sistema}
\gusttekst
1.\quad Rečima – najčešće željeni način ponašanja – behavioral model\\
2.\quad Funkcionalnom tabelom – praktično za sisteme sa malim brojem ulaza i izlaza\\
3.\quad Jedačinama Bulove algebre – kombinacione mreže\\
4.\quad Električnom šemom – pogodno za realizaciju digitalnog sistema, omogućuje jedna vid sinteze i simulacije sistema - Schematic Capture\par
\hspace{5mm}\begin{minipage}{.93\textwidth}
1.\quad Deo CAD Computer Aided Design\\
2.\quad Deo EDA Electronic Design Automation – ECAD\\
3.\quad Najčešće za potrebe pravljenja PCB-a (Printed Circuit Board)
\end{minipage}\par\vspace{4mm}
Danas - za složene i programabilne sisteme mnogo efikasniji\par
\begin{center}HDL – Hardware Description Language\end{center}
Slični tipičnim programskim jezicima – dele neke karakteristike sa njia (objektno orjentisano programiranje, konkuretno programiranje, itd.) ali im je osnovna namena OPIS hardvera u odnosu na posao koji treba neki procesor da izvrši.
\end{frame}
